\chapter*{付録 B：Merton 型構造モデルと B\&S 推定式の数学的関係}
\addcontentsline{toc}{chapter}{付録 B：Merton 型構造モデルと B\&S 推定式の数学的関係}
\label{app:merton_bs}

本付録では，本研究が依拠する理論的基盤として，
Merton（1974）型構造モデルと Bharath and Shumway（2008）（以下 B\&S）の
推定式との数学的関係を体系的に整理する。
構造モデルに基づく倒産距離（Distance to Default; DD）の導出，
B\&S によるロジット推定式との接続，
さらに本研究の潜在変数モデルおよび Ordered Logit への写像を明示することで，
理論的整合性と数学的厳密性を確認する。

% ------------------------------------------------------------
\section*{B.1　Merton 型構造モデルと倒産距離 DD}
\addcontentsline{toc}{section}{B.1　Merton 型構造モデルと倒産距離 DD}

Merton（1974）の構造モデルでは，企業価値 $V_t$ が幾何ブラウン運動に従うと仮定される。



\[
\frac{\mathrm{d}V_t}{V_t}
=
\mu\,\mathrm{d}t + \sigma\,\mathrm{d}W_t,
\]



ここで $\mu$ は期待成長率，$\sigma$ はボラティリティ，
$W_t$ は標準ブラウン運動である。
満期 $T$ において企業価値 $V_T$ が負債額 $F$ を下回るとき，
株主にとって企業は倒産とみなされる。

幾何ブラウン運動の性質より，



\[
\ln V_T
\sim
\mathcal{N}
\left(
\ln V_0 + (\mu - \tfrac{1}{2}\sigma^2)T,\;
\sigma^2 T
\right),
\]



となるため，倒産確率は



\[
\Pr(V_T < F)
=
\Phi\left(
-\frac{\ln(V_0/F) + (\mu - \tfrac{1}{2}\sigma^2)T}{\sigma\sqrt{T}}
\right),
\]



と表される。
分子の項



\[
DD
=
\frac{\ln(V_0/F) + (\mu - \tfrac{1}{2}\sigma^2)T}{\sigma\sqrt{T}}
\]



が倒産距離（Distance to Default）であり，
企業価値が倒産閾値から何標準偏差離れているかを表す。
この DD は確率論に基づき厳密に導出される量であり，
本研究の潜在倒産距離の基礎となる。

% ------------------------------------------------------------
\section*{B.2　B\&S 推定式とロジットモデル}
\addcontentsline{toc}{section}{B.2　B\&S 推定式とロジットモデル}

Bharath and Shumway（2008）は，Merton 型モデルに基づく DD を
実務的に近似しつつ，倒産確率をロジットモデルで推定する枠組みを提示した。
彼らのアプローチは次の二段階構造を持つ。

\begin{enumerate}
  \item 市場データ・財務データから企業価値 $V$ や負債額 $F$ を近似し，DD を計算する。
  \item DD を含む説明変数ベクトル $X$ に基づき，倒産確率をロジットで推定する。
\end{enumerate}

ロジット推定式は概念的に次のように表される。



\[
\Pr(\text{Default}_{i,t+1}=1 \mid X_{i,t})
=
\Lambda\left(
\alpha + \beta_1 \text{DD}_{i,t} + \beta_2 \text{Size}_{i,t}
+ \beta_3 \text{Leverage}_{i,t} + \cdots
\right),
\]



ここで $\Lambda(\cdot)$ はロジスティック分布の累積分布関数である。

重要なのは，B\&S が Merton 型モデルを「厳密な閉形式解」としてではなく，
倒産リスクの構造を捉えるための近似的基盤として用いている点である。
すなわち，DD は構造モデルから導出されるが，
最終的な倒産確率の推定はロジットという離散選択モデルの枠組みで行われる。

この構造は，本研究が採用する
「Merton 型モデル → 潜在変数 → Ordered Logit」
という接続と数学的に同型である。

% ------------------------------------------------------------
\section*{B.3　本研究の潜在変数モデルとの接続}
\addcontentsline{toc}{section}{B.3　本研究の潜在変数モデルとの接続}

本研究では，Merton 型モデルにおける DD の構造を踏まえつつ，
新電力企業の倒産リスクを「市場依存型倒産距離」として再定義する。
第3章で導出した潜在倒産距離 $Z^{\ast}_i$ は，



\[
Z^{\ast}_i
=
a \mu^{\ast}_i - b \sigma^{\ast}_i + c S^{\ast}_i - d D^{\ast}_i,
\]



のような線形構造として定式化される。
これは Merton 型 DD の一次近似として解釈でき，
収益性 $\mu^{\ast}$，市場ショック $\sigma^{\ast}$，
支援能力 $S^{\ast}$，負債比率 $D^{\ast}$ が
Merton 型モデルの $\mu$，$\sigma$，$V/F$ の役割を
産業特性に応じて再パラメータ化したものである。

潜在変数に基づく倒産確率はロジスティック関数により



\[
p_i
=
\frac{1}{1 + \exp(-Z^{\ast}_i)}
\]



と表される。
これは B\&S のロジット推定式と数学的に同型である。

さらに，本研究では倒産を二値ではなく複数段階の現象として扱うため，
潜在変数 $y^*_{i,t}$ を Ordered Logit の枠組みに拡張する。



\[
y^*_{i,t}
=
\gamma_0 + \gamma_1 \mu_{i,t} + \gamma_2 \sigma_{i,t}
+ \gamma_3 S_{i,t} + \gamma_4 D_{i,t} + \varepsilon_{i,t},
\]



潜在変数が複数の閾値 $\tau_k$ を跨ぐことで倒産ステージが決定される。



\[
\text{default\_stage}_{i,t} =
\begin{cases}
0 & (y^*_{i,t} < \tau_1) \\
1 & (\tau_1 \le y^*_{i,t} < \tau_2) \\
2 & (\tau_2 \le y^*_{i,t})
\end{cases}
\]



この枠組みは，
Merton → DD → ロジット  
という構造を  
Merton → 潜在変数 $y^*$ → Ordered Logit  
へと拡張したものであり，
潜在変数モデルとしての数学的厳密性を保持している。

本付録では，Merton 型構造モデルに基づく倒産距離 DD の数学的構造を起点とし，
Bharath \& Shumway（2008）が行った一般化（proxyDD）との対応関係を整理する。
さらに，この一般化された倒産距離を一次テイラー展開により線形化し，
本研究で用いる構造的線形倒産距離 \(Z_{\text{lin}}\) の導出過程を示す。
この導出の全体像を図 \ref{fig:zlin_flow} に示す。

\begin{figure}[htbp]
\centering
\begin{tikzpicture}[node distance=2.2cm]

% --- Merton Model ---
\node[rectangle, draw, rounded corners, text width=9cm, align=center] (merton)
{\textbf{Merton 型倒産距離（DD）}\par
\(\displaystyle 
DD = \frac{\log(V/Thr) + (\mu_V - 0.5\sigma_V^2)T}{\sigma_V\sqrt{T}}
\)
};

% --- B&S Generalization ---
\node[rectangle, draw, rounded corners, text width=9cm, align=center, below=of merton] (bs)
{\textbf{B\&S による構造の抽象化}\par
一般化された倒産距離\par
\(\displaystyle DD = f(\mu, \sigma, S, Thr)\)
};

% --- First-order Taylor ---
\node[rectangle, draw, rounded corners, text width=9cm, align=center, below=of bs] (taylor)
{\textbf{一次テイラー展開（線形近似）}\par
\(\displaystyle 
DD \approx 
\alpha_0 + \alpha_1(\mu-\mu_0)
+ \alpha_2(\sigma-\sigma_0)
+ \alpha_3(S-S_0)
+ \alpha_4(Thr-Thr_0)
\)
};

% --- Reparameterization ---
\node[rectangle, draw, rounded corners, text width=9cm, align=center, below=of taylor] (zlin)
{\textbf{本研究の構造的線形倒産距離}\par
\(\displaystyle 
Z_{\text{lin}} = a\mu - b\sigma + cS - dThr
\)
};

% --- Arrows ---
\draw[->, thick] (merton.south) --
    node[midway, right]{\shortstack[l]{観測不能な企業価値を観測可能な proxy に写像\\
    平均 − 閾値／不確実性の構造を一般化}}
    (bs.north);

\draw[->, thick] (bs.south) --
    node[midway, right]{代表点近傍で一次テイラー展開（線形化）}
    (taylor.north);

\draw[->, thick] (taylor.south) --
    node[midway, right]{非線形 DD を新電力企業の観測可能変数へ線形写像}
    (zlin.north);


\end{tikzpicture}

\caption{Merton 型倒産距離から構造的線形倒産距離 \(Z_{\text{lin}}\) への導出フロー}
\label{fig:zlin_flow}
\end{figure}


図 \ref{fig:zlin_flow} が示すように，Merton 型倒産距離 DD は企業価値の
乗数的変化を対数変換によって線形化した構造を持つ。
B\&S はこの構造を保持したまま観測可能な財務データへ写像し，
一般化された倒産距離 \(DD=f(\mu,\sigma,S,Thr)\) を導出している。
本研究はこの一般化された DD を代表点近傍で一次テイラー展開することで線形化し，
新電力企業の観測可能変数に基づく構造的線形倒産距離
\(Z_{\text{lin}} = a\mu - b\sigma + cS - dThr\) を構築する。

さらに，潜在変数を Ordered Logit の閾値構造へ写像することで，
倒産を二値ではなく複数段階の現象として扱うことが可能となる。
理論編では倒産プロセスを忠実に表現するため 4 段階構造を採用し，
実証編では観測可能性の制約を踏まえて 3 区分に簡略化した。

このように，Merton → B\&S → 潜在変数 → Ordered Logit（4段階／3段階）
という連続的な写像構造により，
本研究は構造モデルの数学的厳密性と実証分析の現実的制約を両立させている。



% ------------------------------------------------------------
\section*{B.4　数学的厳密性と本研究の位置づけ}
\addcontentsline{toc}{section}{B.4　数学的厳密性と本研究の位置づけ}

以上の整理から，本研究の理論的枠組みは次のように位置づけられる。

\begin{itemize}
  \item \textbf{Merton 型構造モデル}\\ 
  幾何ブラウン運動と満期時点での倒産判定に基づき，
  倒産距離 DD を厳密に導出する確率論的基盤。

  \item \textbf{B\&S 推定式}\\ 
  DD を含む説明変数に基づき，
  ロジットモデルで倒産確率を推定する実証的枠組み。
  構造モデルの近似を用いつつ，離散選択モデルとしての数学的整合性を維持する。

  \item \textbf{本研究の潜在変数モデルと Ordered Logit}\\
  Merton 型 DD の構造を新電力企業の文脈に即して再パラメータ化し，
  潜在変数と閾値により倒産ステージを推定する枠組み。
  B\&S 型ロジットを Ordered Logit へ拡張したものであり，
  潜在変数モデルとしての数学的厳密性を保持する。
\end{itemize}

本付録で示したように，本研究は Merton 型構造モデルを
倒産距離 DD の導出と潜在変数モデルの構造において明示的に用いている。
そのうえで，B\&S 型ロジット推定と Ordered Logit を通じて，
新電力企業の倒産リスクを複数段階の現象として推定するという
理論と実証を接続する枠組みを構築している。
この点で，本研究は数学的厳密性と実証的有用性を両立する
構造モデルベースの倒産分析として位置づけられる。
